% !TeX root = ../colloquium-2.tex

% Формулировка правил дифференцирования.

\phantomsection\label{sec_Definition07}
\begin{definition}[Формулировка правил дифференцирования.]
    Пусть функции $f$ и $g$ дифференцируемы в точке $a$. \\
    Тогда:
    \begin{enumerate}
        \item \textbf{Свойство линейности:} \\
        Функция $\alpha f + \beta g$ дифференцируема в точке $a$ при любых 
        $\alpha, \; \beta \in \mathbb{R}$ и 
        $$ (\alpha f + \beta g)' (a) = \alpha f'(a) + \beta g' (a) $$

        \item \textbf{Правило Лейбница:} \\
        $f \cdot g$ дифференцируема в точке $a$, а также $(f \cdot g)'(a) = f'(a)g(a) + f(a)g'(a)$

        \item Если $g \neq 0$ в некоторой окрестности точки $a$, то $\frac{f}{g}$ дифференцируема
        в точках $a$ и $\big( \frac{f}{g} \big)' (a) = \frac{f'(a)g(a) - g'(a)f(a)}{g^2(a)}$.
    \end{enumerate}
\end{definition}